\section{Antiferromagnetic spin waves}\label{app:D.19}

In a collinear antiferromagnet the classical ground state has the spins of one sublattice up and those of the other down. A spin flip on either sublattice is not an eigenstate, because the exchange term $S^+_iS^-_j$ between opposite sublattices creates two flips at once. The spin waves are therefore Bogoliubov mixtures of flips on the two sublattices, and their energy grows linearly at small wavevector instead of quadratically as in a ferromagnet. This section derives the dispersion with the Holstein--Primakoff transformation~\cite{holstein1940field} at linear order, for exchange of any range within and between the sublattices. The result, \cref{eq:D.19-omega}, is used in \cref{sec:7.1} and \cref{app:E.9}; \cref{sec:14.1} computes the spin weights from it.

\subsection*{The two-sublattice Heisenberg model}

The Hamiltonian is the Heisenberg model \cref{eq:7.1-heis}, $H_{\mathrm{H}} = \sum_{\langle ij\rangle}J^{\mathrm{ex}}_{ij}\vect{S}_i\cdot\vect{S}_j$, with each pair counted once and $J^{\mathrm{ex}}_{ij} > 0$ antiferromagnetic. The magnetic ions form two interpenetrating sublattices A and B with $N_{\mathrm{s}}$ sites each, related by a translation or an inversion that exchanges them, and every ion has spin length $|\vect{S}|$. In the classical reference state the spins of A point along $+\hat{\vect{z}}$ and those of B along $-\hat{\vect{z}}$. Exchange of any range enters through two lattice sums, over the neighbors of a site of A on its own sublattice and on the other,
\begin{equation}
  J^{\mathrm{ex}}_{\mathrm{AA}}(\vect{k}) = \sum_{\vect{\delta}\in\mathrm{A},\,\vect{\delta}\ne0}J^{\mathrm{ex}}_{0\vect{\delta}}\,\ee^{\ii\vect{k}\cdot\vect{\delta}}, \qquad
  J^{\mathrm{ex}}_{\mathrm{AB}}(\vect{k}) = \sum_{\vect{\delta}\in\mathrm{B}}J^{\mathrm{ex}}_{0\vect{\delta}}\,\ee^{\ii\vect{k}\cdot\vect{\delta}} ,
  \label{eq:D.19-Jk}
\end{equation}
where $\vect{\delta}$ runs over the positions of the neighbors relative to the site. The symmetry that exchanges A and B gives the same sums for a site of B. When the set of neighbors is symmetric under $\vect{\delta} \to -\vect{\delta}$, as it is for every lattice of this section, both sums are real and even in $\vect{k}$. In the type-II order of NiO, ferromagnetic $(111)$ planes of the fcc nickel lattice alternate in sign (\cref{app:E.9}); all 6 next-nearest neighbors and 6 of the 12 nearest neighbors lie on the other sublattice, and the other 6 nearest neighbors on the same one.

\subsection*{Holstein--Primakoff bosons}

On a site of A, the operators
\begin{equation}
  S^z_i = |\vect{S}| - d^\dagger_{\mathrm{A}i}d_{\mathrm{A}i}, \qquad
  S^+_i = \Bigl(2|\vect{S}| - d^\dagger_{\mathrm{A}i}d_{\mathrm{A}i}\Bigr)^{1/2}d_{\mathrm{A}i}, \qquad
  S^-_i = (S^+_i)^\dagger ,
  \label{eq:D.19-HP}
\end{equation}
obey the spin commutation relations when $d_{\mathrm{A}i}$ is a boson, $[d_{\mathrm{A}i},d^\dagger_{\mathrm{A}i}] = 1$, on the physical states $d^\dagger_{\mathrm{A}i}d_{\mathrm{A}i} \le 2|\vect{S}|$~\cite{holstein1940field}. The boson number counts the deviation from the fully polarized site. On a site of B, where the reference spin points down, the roles of raising and lowering are exchanged: $S^z_j = -|\vect{S}| + d^\dagger_{\mathrm{B}j}d_{\mathrm{B}j}$ and $S^-_j = (2|\vect{S}| - d^\dagger_{\mathrm{B}j}d_{\mathrm{B}j})^{1/2}d_{\mathrm{B}j}$. Linear spin-wave theory keeps the leading order in $1/|\vect{S}|$: it replaces the square roots by $(2|\vect{S}|)^{1/2}$ and drops all terms of the Hamiltonian beyond the second order in the bosons.

For a pair on opposite sublattices, $i \in \mathrm{A}$ and $j \in \mathrm{B}$,
\begin{equation*}
  \vect{S}_i\cdot\vect{S}_j = S^z_iS^z_j + \tfrac12\bigl(S^+_iS^-_j + S^-_iS^+_j\bigr)
  \simeq -|\vect{S}|^2 + |\vect{S}|\bigl(d^\dagger_{\mathrm{A}i}d_{\mathrm{A}i} + d^\dagger_{\mathrm{B}j}d_{\mathrm{B}j}\bigr) + |\vect{S}|\bigl(d_{\mathrm{A}i}d_{\mathrm{B}j} + d^\dagger_{\mathrm{A}i}d^\dagger_{\mathrm{B}j}\bigr) ,
\end{equation*}
since $S^+_iS^-_j \simeq 2|\vect{S}|\,d_{\mathrm{A}i}d_{\mathrm{B}j}$. Each deviation raises the energy of the antiparallel bond, and the transverse term creates or destroys two deviations together. For a pair on the same sublattice, both up,
\begin{equation*}
  \vect{S}_i\cdot\vect{S}_j \simeq |\vect{S}|^2 - |\vect{S}|\bigl(d^\dagger_id_i + d^\dagger_jd_j\bigr) + |\vect{S}|\bigl(d^\dagger_id_j + d^\dagger_jd_i\bigr) ,
\end{equation*}
the ferromagnetic form, in which deviations hop; the same holds on B. The quartic term $d^\dagger d\,d^\dagger d$ of $S^z_iS^z_j$ has been dropped.

\subsection*{Fourier transform}

Sum over pairs. Every site of A has the antiparallel bonds of total strength $J^{\mathrm{ex}}_{\mathrm{AB}}(0)$ and the parallel bonds $J^{\mathrm{ex}}_{\mathrm{AA}}(0)$, and likewise every site of B. Introduce
\begin{equation*}
  d_{\mathrm{A}i} = N_{\mathrm{s}}^{-1/2}\sum_{\vect{k}}\ee^{\ii\vect{k}\cdot\vect{r}_i}\,d_{\mathrm{A}\vect{k}}, \qquad
  d_{\mathrm{B}j} = N_{\mathrm{s}}^{-1/2}\sum_{\vect{k}}\ee^{-\ii\vect{k}\cdot\vect{r}_j}\,d_{\mathrm{B}\vect{k}} ,
\end{equation*}
with $\vect{k}$ in the Brillouin zone of one sublattice. The opposite sign on B is a convenience: $d_{\mathrm{B}\vect{k}}$ destroys a deviation of crystal momentum $-\vect{k}$, so that the pair terms couple $d_{\mathrm{A}\vect{k}}$ with $d_{\mathrm{B}\vect{k}}$. With $\sum_i\ee^{\ii(\vect{k}-\vect{k}')\cdot\vect{r}_i} = N_{\mathrm{s}}\delta_{\vect{k}\vect{k}'}$, the hopping terms on A give $|\vect{S}|\sum_{\vect{k}}J^{\mathrm{ex}}_{\mathrm{AA}}(\vect{k})\,d^\dagger_{\mathrm{A}\vect{k}}d_{\mathrm{A}\vect{k}}$, where each pair counted once contributes $J^{\mathrm{ex}}_{0\vect{\delta}}$ to both $\ee^{\ii\vect{k}\cdot\vect{\delta}}$ and $\ee^{-\ii\vect{k}\cdot\vect{\delta}}$, and those on B give the same with $\vect{k} \to -\vect{k}$, which changes nothing. The pair terms give $\sum_i\sum_{\vect{\delta}\in\mathrm{B}}J^{\mathrm{ex}}_{0\vect{\delta}}\,d_{\mathrm{A}i}d_{\mathrm{B},i+\vect{\delta}} = \sum_{\vect{k}}J^{\mathrm{ex}}_{\mathrm{AB}}(-\vect{k})\,d_{\mathrm{A}\vect{k}}d_{\mathrm{B}\vect{k}}$. Collecting,
\begin{equation}
  \begin{gathered}
  H_{\mathrm{H}} \simeq E_{\mathrm{cl}} + |\vect{S}|\sum_{\vect{k}}\Bigl[A_{\vect{k}}\bigl(d^\dagger_{\mathrm{A}\vect{k}}d_{\mathrm{A}\vect{k}} + d^\dagger_{\mathrm{B}\vect{k}}d_{\mathrm{B}\vect{k}}\bigr) + B_{\vect{k}}\bigl(d_{\mathrm{A}\vect{k}}d_{\mathrm{B}\vect{k}} + d^\dagger_{\mathrm{B}\vect{k}}d^\dagger_{\mathrm{A}\vect{k}}\bigr)\Bigr] , \\
  A_{\vect{k}} = J^{\mathrm{ex}}_{\mathrm{AB}}(0) - J^{\mathrm{ex}}_{\mathrm{AA}}(0) + J^{\mathrm{ex}}_{\mathrm{AA}}(\vect{k}), \qquad B_{\vect{k}} = J^{\mathrm{ex}}_{\mathrm{AB}}(\vect{k}) ,
  \end{gathered}
  \label{eq:D.19-Hk}
\end{equation}
with the classical energy $E_{\mathrm{cl}} = N_{\mathrm{s}}|\vect{S}|^2[J^{\mathrm{ex}}_{\mathrm{AA}}(0) - J^{\mathrm{ex}}_{\mathrm{AB}}(0)]$, which counts each bond once. The coefficient $A_{\vect{k}}$ is the cost of one deviation, the field of the antiparallel neighbors minus that of the parallel ones, plus the ferromagnetic hopping on its own sublattice. For $J^{\mathrm{ex}}_{\mathrm{AB}} = 0$ and $J^{\mathrm{ex}}_{\mathrm{AA}} < 0$, \cref{eq:D.19-Hk} is a ferromagnet with $|\vect{S}|[J^{\mathrm{ex}}_{\mathrm{AA}}(\vect{k}) - J^{\mathrm{ex}}_{\mathrm{AA}}(0)]$ per mode, the one-magnon energy of \cref{sec:7.1}.

\subsection*{Bogoliubov transformation}

At each $\vect{k}$, \cref{eq:D.19-Hk} is the two-mode Hamiltonian \cref{eq:C.7-Hboson} with $d_1 = d_{\mathrm{A}\vect{k}}$, $d_2 = d_{\mathrm{B}\vect{k}}$, the diagonal energy $|\vect{S}|A_{\vect{k}}$ and the pairing amplitude $g_0 = |\vect{S}|B_{\vect{k}}$. \Cref{app:C.7} diagonalizes it with the boson operators
\begin{equation}
  a_{\vect{k}1} = u^{\mathrm{M}}_{\vect{k}}\,d_{\mathrm{A}\vect{k}} + v^{\mathrm{M}}_{\vect{k}}\,d^\dagger_{\mathrm{B}\vect{k}}, \qquad
  a_{\vect{k}2} = u^{\mathrm{M}}_{\vect{k}}\,d_{\mathrm{B}\vect{k}} + v^{\mathrm{M}}_{\vect{k}}\,d^\dagger_{\mathrm{A}\vect{k}} ,
  \label{eq:D.19-bogo}
\end{equation}
with $u^{\mathrm{M}}_{\vect{k}} = \cosh\theta_{\vect{k}}$, $v^{\mathrm{M}}_{\vect{k}} = \sinh\theta_{\vect{k}}$ and $\tanh2\theta_{\vect{k}} = B_{\vect{k}}/A_{\vect{k}}$. Equivalently,
\begin{equation*}
  (u^{\mathrm{M}}_{\vect{k}})^2 = \frac12\Bigl(\frac{A_{\vect{k}}}{\sqrt{A_{\vect{k}}^2 - B_{\vect{k}}^2}} + 1\Bigr), \qquad
  (v^{\mathrm{M}}_{\vect{k}})^2 = \frac12\Bigl(\frac{A_{\vect{k}}}{\sqrt{A_{\vect{k}}^2 - B_{\vect{k}}^2}} - 1\Bigr), \qquad
  u^{\mathrm{M}}_{\vect{k}}v^{\mathrm{M}}_{\vect{k}} = \frac{B_{\vect{k}}}{2\sqrt{A_{\vect{k}}^2 - B_{\vect{k}}^2}} ,
\end{equation*}
from $\cosh2\theta = A/(A^2 - B^2)^{1/2}$ and $\sinh2\theta = B/(A^2 - B^2)^{1/2}$. By \cref{eq:C.7-Eboson},
\begin{equation}
  \begin{gathered}
  H_{\mathrm{H}} \simeq E_{\mathrm{gs}} + \sum_{\vect{k}}\omega^{\mathrm{M}}_{\vect{k}}\bigl(a^\dagger_{\vect{k}1}a_{\vect{k}1} + a^\dagger_{\vect{k}2}a_{\vect{k}2}\bigr), \qquad
  \omega^{\mathrm{M}}_{\vect{k}} = |\vect{S}|\bigl(A_{\vect{k}}^2 - B_{\vect{k}}^2\bigr)^{1/2} , \\
  E_{\mathrm{gs}} = E_{\mathrm{cl}} + \sum_{\vect{k}}\bigl(\omega^{\mathrm{M}}_{\vect{k}} - |\vect{S}|A_{\vect{k}}\bigr) .
  \end{gathered}
  \label{eq:D.19-omega}
\end{equation}
With $A_{\vect{k}}$ and $B_{\vect{k}}$ from \cref{eq:D.19-Hk}, this is \cref{eq:7.1-afm}, and also \cref{eq:E.9-afm}. The two branches are degenerate without anisotropy. They differ in the spin they carry: $a^\dagger_{\vect{k}1} = u^{\mathrm{M}}_{\vect{k}}d^\dagger_{\mathrm{A}\vect{k}} + v^{\mathrm{M}}_{\vect{k}}d_{\mathrm{B}\vect{k}}$ adds a deviation on A or removes one from B, and either lowers the total $S^z$ by one, while $a^\dagger_{\vect{k}2}$ raises it by one. The ground-state energy lies below the classical one by $\sum_{\vect{k}}(|\vect{S}|A_{\vect{k}} - \omega^{\mathrm{M}}_{\vect{k}}) \ge 0$: the Néel state is not the ground state, and the zero-point motion of the spin waves lowers the energy. In the true ground state, annihilated by every $a_{\vect{k}\nu}$, a site of A carries the deviation $\langle d^\dagger_{\mathrm{A}\vect{k}}d_{\mathrm{A}\vect{k}}\rangle = (v^{\mathrm{M}}_{\vect{k}})^2$ per mode, by \cref{app:C.7}; its sum over the zone is the zero-point reduction of the ordered moment that \cref{sec:14.1} evaluates.

The theory is consistent only if $A_{\vect{k}} \ge |B_{\vect{k}}|$ for every $\vect{k}$. Otherwise $\omega^{\mathrm{M}}_{\vect{k}}$ is imaginary, the deviations grow, and the assumed classical order is not a minimum of the classical energy. The expansion is in $1/|\vect{S}|$: the quartic terms that were dropped change the energies by a relative amount of order $1/(2|\vect{S}|)$ times a lattice sum of $(v^{\mathrm{M}}_{\vect{k}})^2$, small for three-dimensional lattices and for $|\vect{S}| \ge 1$.

\subsection*{Linear dispersion}

At $\vect{k} = 0$, $A_0 = B_0 = J^{\mathrm{ex}}_{\mathrm{AB}}(0)$ and $\omega^{\mathrm{M}}_0 = 0$: rotating all spins together costs nothing, the Goldstone mode of the broken spin-rotation symmetry. For small $k$, write $A_{\vect{k}}^2 - B_{\vect{k}}^2 = (A_{\vect{k}} - B_{\vect{k}})(A_{\vect{k}} + B_{\vect{k}})$. The sum tends to $2J^{\mathrm{ex}}_{\mathrm{AB}}(0)$, and the difference is quadratic, $A_{\vect{k}} - B_{\vect{k}} = [J^{\mathrm{ex}}_{\mathrm{AB}}(0) - J^{\mathrm{ex}}_{\mathrm{AB}}(\vect{k})] - [J^{\mathrm{ex}}_{\mathrm{AA}}(0) - J^{\mathrm{ex}}_{\mathrm{AA}}(\vect{k})] \simeq \frac12\sum_{\vect{\delta}\in\mathrm{B}}J^{\mathrm{ex}}_{0\vect{\delta}}(\vect{k}\cdot\vect{\delta})^2 - \frac12\sum_{\vect{\delta}\in\mathrm{A}}J^{\mathrm{ex}}_{0\vect{\delta}}(\vect{k}\cdot\vect{\delta})^2$. Hence
\begin{equation}
  \omega^{\mathrm{M}}_{\vect{k}} \simeq |\vect{S}|\Bigl\{J^{\mathrm{ex}}_{\mathrm{AB}}(0)\Bigl[\sum_{\vect{\delta}\in\mathrm{B}}J^{\mathrm{ex}}_{0\vect{\delta}}(\vect{k}\cdot\vect{\delta})^2 - \sum_{\vect{\delta}\in\mathrm{A}}J^{\mathrm{ex}}_{0\vect{\delta}}(\vect{k}\cdot\vect{\delta})^2\Bigr]\Bigr\}^{1/2} ,
  \label{eq:D.19-linear}
\end{equation}
linear in $k$, with a velocity that is the geometric mean of the exchange field $|\vect{S}|J^{\mathrm{ex}}_{\mathrm{AB}}(0)$ and the stiffness term in brackets. A ferromagnet has no pair term and $\omega \propto k^2$; the antiferromagnet's linear law comes from the mixing of deviations on the two sublattices, which costs the large exchange field.

For the type-II order of the fcc lattice with only the next-nearest-neighbor exchange $J_2$, the 6 antiparallel neighbors sit at $\pm a\hat{\vect{x}}$, $\pm a\hat{\vect{y}}$, $\pm a\hat{\vect{z}}$ with $a$ the cubic lattice constant, $J^{\mathrm{ex}}_{\mathrm{AB}}(\vect{k}) = 2J_2(\cos k_xa + \cos k_ya + \cos k_za)$ and $J^{\mathrm{ex}}_{\mathrm{AA}} = 0$. Then $\omega^{\mathrm{M}}_{\vect{k}} = 6J_2|\vect{S}|[1 - \zeta_{\vect{k}}^2]^{1/2}$ with $\zeta_{\vect{k}} = \frac13\sum_\alpha\cos k_\alpha a$. The maximum, $6J_2|\vect{S}|$, is reached where $\zeta_{\vect{k}} = 0$, and the small-$k$ law \cref{eq:D.19-linear}, with $\sum_{\vect{\delta}}(\vect{k}\cdot\vect{\delta})^2 = 2a^2k^2$, is $\omega^{\mathrm{M}}_{\vect{k}} \simeq 2\sqrt3\,J_2|\vect{S}|\,ak$. The nearest-neighbor exchange $J_1$ adds 6 antiparallel bonds to $J^{\mathrm{ex}}_{\mathrm{AB}}$ and 6 parallel bonds to $J^{\mathrm{ex}}_{\mathrm{AA}}$ and changes these numbers at first order in $J_1/J_2$.

\subsection*{The spin operators}

A probe that couples to the electron spin density, as in \cref{sec:7.1}, needs the spin operators expressed through the magnon operators. With the Fourier sign of \cref{sec:7.1}, $S^\pm(\vect{k}) = \sum_j\ee^{\ii\vect{k}\cdot\vect{r}_j}S^\pm_j$, the part on A is to linear order $S^+_{\mathrm{A}}(\vect{k}) \equiv \sum_{i\in\mathrm{A}}\ee^{\ii\vect{k}\cdot\vect{r}_i}S^+_i = (2|\vect{S}|N_{\mathrm{s}})^{1/2}d_{\mathrm{A},-\vect{k}}$, by the lattice sum $\sum_i\ee^{\ii(\vect{k}+\vect{k}')\cdot\vect{r}_i} = N_{\mathrm{s}}\delta_{\vect{k},-\vect{k}'}$. Inverting \cref{eq:D.19-bogo} as in \cref{app:C.7},
\begin{equation}
  d_{\mathrm{A}\vect{k}} = u^{\mathrm{M}}_{\vect{k}}\,a_{\vect{k}1} - v^{\mathrm{M}}_{\vect{k}}\,a^\dagger_{\vect{k}2}, \qquad
  d_{\mathrm{B}\vect{k}} = u^{\mathrm{M}}_{\vect{k}}\,a_{\vect{k}2} - v^{\mathrm{M}}_{\vect{k}}\,a^\dagger_{\vect{k}1} .
  \label{eq:D.19-inverse}
\end{equation}
A transverse spin operator acting on the ground state therefore creates one magnon, with an amplitude $(2|\vect{S}|N_{\mathrm{s}})^{1/2}$ times $u^{\mathrm{M}}_{\vect{k}}$ or $v^{\mathrm{M}}_{\vect{k}}$. On B the raising operator is $S^+_j \simeq (2|\vect{S}|)^{1/2}d^\dagger_{\mathrm{B}j}$, and the same lattice sum turns $\sum_{j\in\mathrm{B}}\ee^{\ii\vect{k}\cdot\vect{r}_j}d^\dagger_{\mathrm{B}j}$ into $N_{\mathrm{s}}^{1/2}d^\dagger_{\mathrm{B},-\vect{k}}$. Since $A_{\vect{k}}$ and $B_{\vect{k}}$, and with them $u^{\mathrm{M}}_{\vect{k}}$ and $v^{\mathrm{M}}_{\vect{k}}$, are even in $\vect{k}$, the sum of both sublattices at small $\vect{k}$ is $S^+(\vect{k}) \simeq (2|\vect{S}|N_{\mathrm{s}})^{1/2}(d_{\mathrm{A},-\vect{k}} + d^\dagger_{\mathrm{B},-\vect{k}}) = (2|\vect{S}|N_{\mathrm{s}})^{1/2}(u^{\mathrm{M}}_{\vect{k}} - v^{\mathrm{M}}_{\vect{k}})(a_{-\vect{k}1} + a^\dagger_{-\vect{k}2})$. Near $\vect{k} = 0$ both coefficients diverge as $(|\vect{S}|A_0/2\omega^{\mathrm{M}}_{\vect{k}})^{1/2}$, but $(u^{\mathrm{M}}_{\vect{k}} - v^{\mathrm{M}}_{\vect{k}})^2 = [(A_{\vect{k}} - B_{\vect{k}})/(A_{\vect{k}} + B_{\vect{k}})]^{1/2} \to 0$, so the response to a uniform transverse field is suppressed. Near the antiferromagnetic wavevector, where the phases of the two sublattices differ by $\pi$, the combination is $u^{\mathrm{M}}_{\vect{k}} + v^{\mathrm{M}}_{\vect{k}}$ and the response is enhanced. The longitudinal operator $S^z$ is quadratic in the bosons and creates magnons in pairs. \Cref{sec:14.1} assembles these weights.

The classical and the linear quantum theories agree at this order. On a periodic fcc cluster of 256 sites with type-II order and $J_1 = 0.1J_2$, the frequencies of the classical precession equations $\dot{\vect{S}}_i = -\vect{S}_i\times\sum_jJ^{\mathrm{ex}}_{ij}\vect{S}_j$, linearized about the Néel state, coincide with \cref{eq:D.19-omega} at all wavevectors of the cluster.
